\documentclass[11pt]{article}
\usepackage[margin=1in]{geometry}
\usepackage{amsmath,amssymb,amsthm}
\usepackage[T1]{fontenc}
\usepackage{lmodern}
\usepackage[hidelinks]{hyperref}
\hypersetup{pdftitle={An explicit joint Laplace transform for Brownian first-visit cell lengths on a circle},pdfauthor={Alec Kriebel}}
\newtheorem{theorem}{Theorem}
\newtheorem{lemma}{Lemma}
\newcommand{\T}{\mathbb T}
\newcommand{\R}{\mathbb R}
\newcommand{\E}{\mathbb E}
\newcommand{\Prob}{\mathbb P}
\newcommand{\ind}{\mathbf 1}
\title{An explicit joint Laplace transform\\for Brownian first-visit cell lengths on a circle}
\author{Alec Kriebel\\\small Independent researcher\\\small\href{https://orcid.org/0009-0001-9320-500X}{ORCID: 0009-0001-9320-500X}}
\date{October 7, 2026}
\begin{document}
\maketitle
\begin{abstract}
We give a complete analytic-transform answer to Georgakopoulos's
question about the lengths of the sets first visited by independent
Brownian particles on a circle, with uniform random or equally spaced
starting points. For every finite number of particles, the joint
Laplace transform is an absolutely convergent series whose coefficients
are finite sums of ordinary integrals of explicit interval-exit kernels.
The integration regions are specified linear inequalities in cumulative
physical hitting times. The outer series has a uniform factorial
remainder bound and uniquely determines the probability measure.
The result also applies to arbitrary fixed distinct starting points.
The construction uses classical heat kernels, strong Markov iteration,
and bounded-support moments. Earlier competing-walk painting results
are credited and distinguished. No compact density, efficient numerical
algorithm, or absolute historical priority is asserted.
\end{abstract}

\section{The question and the model}
Georgakopoulos's question on printed p.~1452 of Oberwolfach Report
23/2018 asks for the joint lengths of first-visit cells of $k$ independent
Brownian particles started uniformly at random or at equally spaced
points of $S^1$ \cite{Georgakopoulos}. The report belongs to the 2018
volume and was published on April 12, 2019. We answer the finite-$k$
distribution question by an explicit transform law; no named density or
large-$k$ limit is imposed in that question.

Normalize the circle to $\T=\R/\mathbb Z$, with Lebesgue measure one.
Let $k\geq1$ and let $W_i$ be independent Brownian motions with generator
$\tfrac12 d^2/dx^2$, starting at distinct labeled points $p_i$.
The particles continue moving independently through all encounters,
seeds, and previously visited territory. Put
\begin{equation}\label{lengths}
 T_i(x)=\inf\{t\geq0:W_i(t)=x\},\qquad
 I(x)=\min\operatorname*{arg\,min}_{1\leq i\leq k}T_i(x),\qquad
 L_i=\int_{\T}\ind_{\{I(x)=i\}}\,dx.
\end{equation}
We justify measurability and the irrelevance of ties below.
For circumference $c>0$, the lengths are $cL_i$, so replace each
Laplace variable $\theta_i$ by $c\theta_i$. A common positive diffusivity
change rescales all clocks equally and leaves ownership unchanged.

\section{Explicit deterministic coefficients}
For $0<\alpha<\ell$ and $t>0$, define
\begin{align}
 f_0(t;\alpha,\ell)
 &=\frac{\pi}{\ell^2}\sum_{n=1}^{\infty}
 n\sin\frac{n\pi\alpha}{\ell}
       e^{-n^2\pi^2t/(2\ell^2)},\label{flux0}\\
 f_1(t;\alpha,\ell)
 &=\frac{\pi}{\ell^2}\sum_{n=1}^{\infty}
 (-1)^{n+1}n\sin\frac{n\pi\alpha}{\ell}
       e^{-n^2\pi^2t/(2\ell^2)}.\label{flux1}
\end{align}
These are the nonnegative densities of interval exit at $0$ and $\ell$,
respectively. The series are absolutely convergent at each $t>0$.
They are summed before any improper integration near $t=0$; no
termwise integration of the signed Fourier series is required.

For a nonempty finite set $A\subset\T$ and $z\notin A$, lift the component
of $\T\setminus A$ containing $z$ to $(a,b)\subset\R$, and represent
$z$ by $\widetilde z\in(a,b)$. Set $\ell=b-a$, $\alpha=\widetilde z-a$,
and, for $q\in A$, define
\begin{equation}\label{kernel}
 K_A(t;z,q)=\ind_{\{a\bmod1=q\}}f_0(t;\alpha,\ell)
            +\ind_{\{b\bmod1=q\}}f_1(t;\alpha,\ell).
\end{equation}
Integer changes of lift do not change this formula. For a singleton
$A=\{q\}$, both endpoints represent $q$, $\ell=1$, and the kernel is
$f_0+f_1$. Previously visited targets are not absorbing.

Fix $m\geq1$ distinct query points $x=(x_1,\ldots,x_m)$, none a seed.
Write $S_m$ for the permutations of $\{1,\ldots,m\}$. For $\pi\in S_m$
and particle $i$, let
\[
 A_r^\pi=\{x_{\pi(r)},\ldots,x_{\pi(m)}\},\qquad
 z_{i,0}=p_i,\quad z_{i,r-1}=x_{\pi(r-1)}\quad(r\geq2).
\]
For $u_i=(u_{i1},\ldots,u_{im})\in(0,\infty)^m$, set
\begin{equation}\label{density}
 D_i^\pi(u_i;x,p_i)=\prod_{r=1}^m
 K_{A_r^\pi}(u_{ir};z_{i,r-1},x_{\pi(r)}).
\end{equation}
For $\boldsymbol\pi=(\pi_1,\ldots,\pi_k)\in S_m^k$, reconstruct the
physical hitting times, rather than comparing the separate increments:
\begin{equation}\label{clock}
 h_{ij}(u,\boldsymbol\pi)=\sum_{r=1}^{\pi_i^{-1}(j)}u_{ir}.
\end{equation}
For a label vector $a=(a_1,\ldots,a_m)\in\{1,\ldots,k\}^m$, put
\begin{align}
 \mathcal C_a(\boldsymbol\pi)
 &=\{u>0:h_{a_j,j}(u,\boldsymbol\pi)<h_{ij}(u,\boldsymbol\pi)
               \text{ for every }j\text{ and }i\ne a_j\},\label{cone}\\
 Q_a(x;p)
 &=\sum_{\boldsymbol\pi\in S_m^k}
   \int_{\mathcal C_a(\boldsymbol\pi)}
   \prod_{i=1}^kD_i^{\pi_i}(u_i;x,p_i)
   \prod_{i=1}^k\prod_{r=1}^mdu_{ir}.\label{ownership}
\end{align}
Thus $Q_a$ is specified by $(m!)^k$ nonnegative, $km$-dimensional
integrals of the explicit scalar series, over strict linear inequalities
with integer coefficients. Impossible orders have zero factors.

For $\theta_i\geq0$, let $\theta_*=\max_i\theta_i$, define
$M_0(\theta;p)=1$, and put
\begin{equation}\label{moments}
 M_m(\theta;p)=\int_{\T^m}
 \sum_{a\in\{1,\ldots,k\}^m}
 \left(\prod_{j=1}^m\theta_{a_j}\right)Q_a(x;p)
 \,dx_1\cdots dx_m\qquad(m\geq1).
\end{equation}
Assign arbitrary bounded values on the excluded query configurations;
they have zero Lebesgue measure. All coefficient definitions are
deterministic. In particular, they contain no unknown Brownian hitting
probability, conditional law, or integral-equation solution.

\begin{theorem}\label{main}
For every finite $k\geq1$ and every distinct seed vector $p$,
$L=(L_1,\ldots,L_k)$ is well defined almost surely, lies in the simplex
$\{l_i\geq0:\sum_i l_i=1\}$, and has the joint Laplace transform
\begin{equation}\label{transform}
 \mathcal L_p(\theta)=\E_p e^{-\sum_i\theta_iL_i}
       =\sum_{m=0}^{\infty}\frac{(-1)^m}{m!}M_m(\theta;p).
\end{equation}
The series is absolutely convergent, $0\leq M_m\leq\theta_*^m$, and,
for every $N\geq0$,
\begin{equation}\label{remainder}
 \left|\mathcal L_p(\theta)-\sum_{m=0}^{N}
                   \frac{(-1)^m}{m!}M_m(\theta;p)\right|
       \leq\frac{\theta_*^{N+1}}{(N+1)!}.
\end{equation}
The coefficients uniquely determine the entire joint probability
measure. For equally spaced seeds use $p_i=(i-1)/k$. For independent
uniform seeds, the answer is
\begin{equation}\label{uniform}
 \mathcal L_{\mathrm{unif}}(\theta)
 =\sum_{m=0}^{\infty}\frac{(-1)^m}{m!}
                    \int_{\T^k}M_m(\theta;p)\,dp,
\end{equation}
with the same bound \eqref{remainder}.
\end{theorem}

\section{Proof}
\begin{lemma}\label{exit}
The kernels in \eqref{flux0}--\eqref{kernel} are the next-target
exit-time/place densities, and
\begin{equation}\label{normalization}
 \int_0^\infty f_0\,dt=1-\alpha/\ell,\qquad
 \int_0^\infty f_1\,dt=\alpha/\ell,\qquad
 \sum_{q\in A}\int_0^\infty K_A(t;z,q)\,dt=1.
\end{equation}
\end{lemma}
\begin{proof}
The classical Dirichlet heat kernel, obtained from the interval sine
basis and space--time scaling, is
\begin{equation}\label{heat}
 H(t;\alpha,\beta)=\frac2\ell\sum_{n\geq1}
 \sin\frac{n\pi\alpha}{\ell}\sin\frac{n\pi\beta}{\ell}
 e^{-n^2\pi^2t/(2\ell^2)}.
\end{equation}
See also Lalley's Exercise~6(C) \cite{Lalley} for the unit interval.
If $\tau$ is the interval exit time, stopping $B_t^2-t$ at
$\tau\wedge t$ bounds $\E(\tau\wedge t)$ uniformly, since the stopped
position stays in $[0,\ell]$. Hence $\tau<\infty$ almost surely.
Optional stopping for the bounded stopped coordinate gives exit
probabilities $1-\alpha/\ell$ and $\alpha/\ell$.

The Markov property therefore gives
\[
 \Prob_\alpha(\tau>t,B_\tau=0)
    =\int_0^\ell H(t;\alpha,\beta)(1-\beta/\ell)\,d\beta.
\]
For $t>0$, locally uniform convergence permits differentiation.
Use $H_t=H_{\beta\beta}/2$, integrate twice by parts, and use the
Dirichlet boundary values. The negative derivative is
$\tfrac12H_\beta(t;\alpha,0)=f_0(t;\alpha,\ell)$.
With weight $\beta/\ell$ it is
$-\tfrac12H_\beta(t;\alpha,\ell)=f_1(t;\alpha,\ell)$.
This proves nonnegativity, the density interpretation, and the first
two identities in \eqref{normalization}, without improper Fourier
interchanges. Circle continuity identifies the next target with the
exit endpoint of its lifted component. Summing both endpoints,
including their identification for a singleton, proves the last identity.
\end{proof}

For comparison, the familiar interval-exit Laplace transforms are
\begin{equation}\label{exitlaplace}
 \int_0^\infty e^{-st}(f_0,f_1)\,dt
 =\left(\frac{\sinh((\ell-\alpha)\sqrt{2s})}
                   {\sinh(\ell\sqrt{2s})},
        \frac{\sinh(\alpha\sqrt{2s})}
                   {\sinh(\ell\sqrt{2s})}\right)\quad(s>0).
\end{equation}
These follow from $v''/2=sv$ with the corresponding endpoint data;
see, for example, Ernst--Shepp \cite{ErnstShepp}. They are normalization
checks, not new interval formulas.

\begin{lemma}\label{finite}
$D_i^\pi$ is the joint density of the order $\pi$ and elapsed
increments between successive newly visited queries for particle $i$.
Moreover,
\[
 \sum_{\pi\in S_m}\int_{(0,\infty)^m}D_i^\pi(u_i)\,du_i=1,
 \quad Q_a(x;p)=\Prob_p(I(x_1)=a_1,\ldots,I(x_m)=a_m),
 \quad \sum_aQ_a(x;p)=1.
\]
\end{lemma}
\begin{proof}
Stop a single walker when it first reaches the finite query set, remove
that reached target, and repeat with the remaining set. At each step
the current point lies outside the remaining targets. These are finite
stopping times by Lemma~\ref{exit}. Distinct queries cannot be reached
simultaneously by a continuous path. The strong Markov property
\cite[Theorem~2.16]{MoertersPeres} gives exactly the product
\eqref{density}; repeated normalization in \eqref{normalization} gives
its total mass. This argument tracks the same physical walker, which
may cross already visited queries and other particles' territory.

Independence gives the product of the $k$ ordered densities.
Equation~\eqref{clock} reconstructs each physical first-hitting time,
so the owner event is exactly \eqref{cone}. A fixed query's individual
hitting-time distribution is atomless by the singleton kernel.
Independent walkers thus have no tie at a fixed query, almost surely.
The finite union over queries is still null; equivalently the equality
boundaries are proper hyperplanes of zero measure for these densities.
Integrating and summing all orders proves the assertions about $Q_a$.
\end{proof}

\begin{proof}[Proof of Theorem~\ref{main}]
On continuous-path space, the visited set up to a fixed time $t$ is
compact. Its distance from $x$ is the infimum of the distances from
positions at rational times in $[0,t]$, together with the position at
$t$. Hence $(\text{path},x)\mapsto T_i(x)$ is jointly measurable.
For each nonseed $x$, every particle reaches it in finite time and the
independent hitting times are atomless. Fubini shows that, almost
surely, locations with any infinite hitting time or an earliest tie
have measure zero. Thus \eqref{lengths} defines a measurable partition
up to a null set and $\sum_iL_i=1$. This is not an assertion that ties
are absent simultaneously at every location.

Put $Z=\sum_i\theta_iL_i=\int_\T\theta_{I(x)}\,dx$, so
$0\leq Z\leq\theta_*$. Tonelli and Lemma~\ref{finite} yield
$\E_p Z^m=M_m(\theta;p)$ and $0\leq M_m\leq\theta_*^m$.
The exponential series is absolutely dominated by $e^{\theta_*}$;
termwise expectation proves \eqref{transform}. Taylor's theorem for
$e^{-z}$ on $z\geq0$ bounds the degree-$N$ remainder by
$z^{N+1}/(N+1)!$, giving \eqref{remainder}.

The homogeneous moment polynomials give all mixed moments, or these
can be recovered by differentiation of the transform at zero, justified
by bounded support. Polynomials are dense in continuous functions on
the compact simplex by Stone--Weierstrass. Consequently its finite
Borel measures are uniquely determined by these moments.
The remainder and absolute bounds are uniform in distinct $p$.
Coincident independent-uniform seeds have measure zero, so integration
in $p$ and interchange with the series prove \eqref{uniform}.
Equally spaced starts are already a fixed-$p$ instance.
\end{proof}

A random rotation of equally spaced seeds leaves lengths unchanged.
If seeds instead receive labels after circular ordering, integrate the
same fixed-$p$ transform against that ordered seed law. Formula
\eqref{uniform} uses independent uniform labels.
For $k=1$ the transform is $e^{-\theta_1}$. If every $\theta_i=t$,
then $M_m=t^m$ and the transform is $e^{-t}$. Uniform starts are
exchangeable; equally spaced starts have cyclic and dihedral symmetry.
Both give $\E L_i=1/k$; full permutation exchangeability for equally
spaced labels is not asserted.

\section{Attribution, priority and verification}
First-arrival painting is an established model. Gomes J\'unior et al.
\cite{Gomes} treat a two-walker periodic lattice and a one-point
scaling probability. Miller \cite{Miller} studies ownership correlations
and asymptotic volume variance on transient graph families.
Chatterjee--Nabahi--Terlov \cite{Chatterjee} analyze expected interfaces
on a discrete cycle. Baccara's joint real-line competitive-range limit
\cite[Section~6, Proposition~8]{Baccara} remains a Brownian functional
characterization; Si \cite{Si} studies strategic painting payoffs,
including a reflected-Brownian interval limit. These are relevant
predecessors with different observables or regimes.

Moment-based transforms and exit-increment constructions are classical;
see Kac's product-moment framework as reviewed by Fitzsimmons--Pitman
\cite{FitzsimmonsPitman} and the multitime single-walker span laws of
R\'egnier et al. \cite{Regnier}. The reduction here makes every
finite-query ownership probability explicit while preserving each
walker's correlated first-hitting times, then determines the joint
circle volume law. No novelty is claimed for the model, interval
kernels, strong Markov property, moment expansion, or moment determinacy.

Our bounded primary-source audit of October 7, 2026 did not identify
an earlier complete characterization of this finite-$k$ circle law,
or a general explicit result directly supplying all its coefficients.
The full Gomes 1996 article and Dicker's related higher-dimensional
thesis \cite{Dicker} were inaccessible; only the former's publisher abstract and
later primary discussions, and the latter's primary index, were
available. Direct Zenodo registry searches also failed. These
limitations, search records, and theorem-scope comparisons are retained
in the supplement. No absolute priority, exhaustive absence, or
independent-discovery claim is made.

The supplied diagnostics reproduce exact finite-chain order, clock,
and moment identities and compare scalar exit fluxes by independent
spectral and reflected-Gaussian formulas. They test supporting
calculations; they are not a Brownian limit proof or a certificate of
numerical quadrature. The proof above establishes the Brownian result.
The outer bound \eqref{remainder} does not bound numerical errors in
the coefficient integrals or a truncation of their inner kernel series.
A compact density, named distribution, and efficient evaluator remain
possible refinements.

\paragraph{AI and review disclosure.}
AI tools were used extensively in developing, drafting, reproducing,
and adversarially verifying this work. Independent AI reviews are not
conventional human peer review. This preprint is unrefereed and has not
undergone conventional human peer review.

\begin{thebibliography}{99}
\bibitem{Georgakopoulos}
A. Georgakopoulos, Voronoi-like decomposition of $S^1$ with randomness,
in \emph{Enumerative Combinatorics}, Oberwolfach Report 23/2018,
\emph{Oberwolfach Reports} 15 (2018), no.~2, 1381--1464, item on p.~1452.
Published April 12, 2019. \href{https://doi.org/10.4171/OWR/2018/23}{doi:10.4171/OWR/2018/23}.
\bibitem{Lalley}
S. Lalley, \emph{Brownian Motion}, course notes, Exercise~6(A)--(C).
\url{https://galton.uchicago.edu/~lalley/Courses/312/BrownianMotion312.pdf}.
\bibitem{MoertersPeres}
P. M\"orters and Y. Peres, \emph{Brownian Motion}, Cambridge University
Press, 2010, Theorem~2.16, pp.~43--44.
\url{https://www.mi.uni-koeln.de/~moerters/book/book.pdf}.
\bibitem{ErnstShepp}
P. Ernst and L. Shepp, On the time for Brownian motion to visit every
point on a circle, \emph{J. Statist. Plann. Inference} 171 (2016), 130--134.
\href{https://doi.org/10.1016/j.jspi.2015.10.010}{doi:10.1016/j.jspi.2015.10.010}.
\bibitem{Gomes}
S.R. Gomes J\'unior, L.S. Lucena, L.R. da Silva and H.J. Hilhorst,
Coloring of a one-dimensional lattice by two independent random walkers,
\emph{Physica A} 225 (1996), 81--88.
\href{https://doi.org/10.1016/0378-4371(95)00424-6}{doi:10.1016/0378-4371(95)00424-6}.
\bibitem{Miller}
J. Miller, Painting a graph with competing random walks,
\emph{Ann. Probab.} 41 (2013), 636--670.
\href{https://doi.org/10.1214/11-AOP713}{doi:10.1214/11-AOP713}.
\bibitem{Chatterjee}
S. Chatterjee, N. Nabahi and G. Terlov, Interface between competing
random walks on a cycle, 2026. \url{https://arxiv.org/abs/2608.21312v1}.
\bibitem{Baccara}
M. Baccara, On the range of competing random walks, 2026.
\url{https://arxiv.org/abs/2607.02406v1}.
\bibitem{Si}
I.M. Si, Strategic geometry of competing first-passage random walks, 2026.
\url{https://arxiv.org/abs/2608.24837v1}.
\bibitem{Dicker}
L. Dicker, Coloring a $d\geq3$ dimensional lattice with two independent
random walks, master's thesis, 2006. Title and year from the primary
UPenn thesis index; full text was unavailable during the audit.
\url{https://www2.math.upenn.edu/~pemantle/papers/Papers.html}.
\bibitem{FitzsimmonsPitman}
P.J. Fitzsimmons and J. Pitman, Kac's moment formula and the Feynman--Kac
formula for additive functionals of a Markov process,
\emph{Stochastic Process. Appl.} 79 (1999), 117--134.
\href{https://doi.org/10.1016/S0304-4149(98)00081-7}{doi:10.1016/S0304-4149(98)00081-7}.
\bibitem{Regnier}
L. R\'egnier, M. Dolgushev, S. Redner and O. B\'enichou,
Complete visitation statistics of one-dimensional random walks,
\emph{Phys. Rev. E} 105 (2022), 064104.
\href{https://doi.org/10.1103/PhysRevE.105.064104}{doi:\nolinkurl{10.1103/PhysRevE.105.064104}}.
\end{thebibliography}
\end{document}
